%%%%%%%%%%%%%%%%%%%%%%%%%%%%%%%%%%%%%%%%%%%%%%%%%%%%%%%%%%%%%%%%%%%%%%%%%%%%%
%% Author: Joao Paulo de Souza Medeiros                                    %%
%%%%%%%%%%%%%%%%%%%%%%%%%%%%%%%%%%%%%%%%%%%%%%%%%%%%%%%%%%%%%%%%%%%%%%%%%%%%%

\documentclass[11pt,a4paper]{article}

\usepackage[top=3cm,bottom=6cm,left=2cm,right=2cm]{geometry} % Assinaturas
\usepackage[utf8]{inputenc}
\usepackage[brazil]{babel}
\usepackage{graphicx}
\usepackage{paralist}
\usepackage{lineno}
\usepackage{color}
\usepackage{indentfirst}
\usepackage{lastpage}
\usepackage{fancyhdr}
\usepackage{lipsum}

\setlength{\headheight}{14pt}
\pagestyle{fancy}
\chead{Primeira Reunião do Grupo do Projeto NP, 1 de abril de 1985}
\cfoot{Página \thepage\ de \pageref{LastPage}}

\renewcommand{\headrulewidth}{0.4pt}% Default \headrulewidth is 0.4pt
\renewcommand{\footrulewidth}{0.4pt}% Default \footrulewidth is 0.4pt

\begin{document}

% ========================================================================= %
\noindent{%
\parbox{0.25\textwidth}%
{\includegraphics[width=0.25\textwidth]{image/ufrn.eps}}%
\parbox{0.65\textwidth}%
{
  \begin{center}%
  \small\textbf{Universidade Federal do Rio Grande do Norte\\
    Centro de Ensino Superior do Seridó\\
    Laboratório de Elementos do Processamento da Informação\\
    Reunião do Grupo do Projeto NP}
  \end{center}
}%
\parbox{0.1\textwidth}%
{\includegraphics[width=0.1\textwidth]{image/labepi-square.eps}}

\linenumbers
\def\linenumberfont{\normalfont\small\sffamily}

\vspace{1em}

% Se houverem reuniões periódicas, deve-se usar o termo `Reunião Ordinária'
% para aquelas agendadas e `Reunião Extraordinária' para aquelas marcadas em
% datas não programadas.
% Se as reuniões são feitas sob demanda, usa-se apenas o termo `Reunião'.
% ------------------------------------------------------------------------- %
\noindent%
{\bf Primeira Reunião} ({\bf Data}: 1 de abril de 1985)\\
{\bf Presentes}:
\begin{inparaenum}[(i)]\small
\item Professor Nome Completo do Docente A,
\item Professor Nome Completo do Docente B,
\item Servidor Nome Completo do Funcionário A,
\item Servidor Nome Completo do Funcionário B,
\item Discente Nome Completo do Aluno A,
\item Discente Nome Completo do Aluno B.
\end{inparaenum}

\vspace{1em}

% ------------------------------------------------------------------------- %
{\bf Informações}:
Descreve-se informações que contextalizam o exposto na próxima seção,
denominada `Decisões'.
Como no exemplo que se segue.
Para a primeira reunião as seguintes tarefas foram realizadas:
\begin{inparaenum}
\item primeira tarefa,
\item segunda tarefa e
\item terceira tarefa;
\end{inparaenum}
que foram apresentadas pelos membros Professor Docente A, Servidor
Funcionário A e pelo Discente Aluno A, respectivamente. 

\lipsum[1]
\lipsum[3]

\vspace{1em}

% ------------------------------------------------------------------------- %
{\bf Decisões}:
Às dez horas do dia 1 de abril de 1985 teve início a primeira reunião do
grupo Nome do Grupo (NG), responsável pelo projeto Nome do Projeto (NP) do
Laboratório de Elementos do Processamento da Informação (LabEPI) do Centro
Superior de Ensino Superior do Seridó (CERES) da Universidade Federal do Rio
Grande do Norte (UFRN).

Nesta primeira reunião ficou decidido entre os membros que a meta de cada
uma das próximas reuniões seria a definição detalhada das seguintas etapas do
projeto:
\begin{compactenum}
\item primeira etapa,
\item segunda etapa,
\item terceira etapa e 
\item última etapa.
\end{compactenum}
Além disso foram discutidos quais equipamentos deveriam ser adquiridos para
que o desenvolvimento do projeto fosse possível.

\lipsum[2]
\lipsum[4]

Às onze horas e trinta minutos, nada mais havendo a tratar, os membros do NG
deram por por encerrada a reunião, de cujos trabalhos, eu, Discente Aluno B,
lavrei a presente ata que, após lida e aprovada, segue devidamente assinada
por todos os membros presentes.

\end{document}

%%%%%%%%%%%%%%%%%%%%%%%%%%%%%%%%%%%%%%%%%%%%%%%%%%%%%%%%%%%%%%%%%%%%%%%%%%%%%
